% Pista ana-26j2-q1 · Anàlisi
% Procedència: PAU Catalunya 2026 · Convocatòria de juny · Sèrie 5 · Exercici 1
\documentclass[11pt,a4paper]{article}
\usepackage{pista}
\pistainfo{Catalunya 2026}{Convocatòria de juny}{Sèrie 5}

\begin{document}

\section*{\textcolor{blauPista}{Exercici 1}}

\noindent Considereu la paràbola $y = f(x)$, amb $f(x) = -x^{2} + 5x$, i la hipèrbola $y = g(x)$, amb $g(x) = 8 - \dfrac{8}{x + 1}$.

\subsection*{1.1. \normalfont En quants punts es tallen les gràfiques d'aquestes dues funcions? Calculeu-los tots. \hfill\textnormal{[0,75~punts]}}

\pista{Planteja l'equació $f(x) = g(x)$ i multiplica ambdós costats per $(x + 1)$ per eliminar el denominador (compte: $x = -1$ queda fora del domini de $g$, però comprovaràs que no és cap de les solucions que trobis). T'ha de quedar una equació cúbica; treu-hi factor comú $x$ per reduir-la a una equació de segon grau, més fàcil de resoldre.}

\subsection*{1.2. \normalfont Calculeu l'àrea de la regió situada entre les dues corbes des de $x = 1$ fins a $x = 3$. \hfill\textnormal{[1~punt]}}

\pista{Primer necessites saber quina de les dues funcions està per damunt de l'altra en aquest interval (per exemple, avaluant-les totes dues en un punt intermedi, com $x = 2$). L'àrea que et demanen és la integral definida de la diferència --la de dalt menys la de baix-- entre $x = 1$ i $x = 3$; per integrar el terme $\dfrac{8}{x + 1}$, recorda que la seva primitiva és $8\ln|x + 1|$.}

\subsection*{1.3. \normalfont Comproveu que $f(x)$ té les arrels als punts d'abscissa $x = 0$ i $x = 5$, i que l'ordenada del vèrtex de la paràbola $y = f(x)$ és $y = 6{,}25$. Justifiqueu que $y = f(x)$ és l'única paràbola que existeix amb aquestes dues propietats. \hfill\textnormal{[0,75~punts]}}

\pista{Les dues primeres comprovacions són directes: avalua $f$ en $x = 0$ i $x = 5$, i troba l'abscissa del vèrtex amb la fórmula $x = -b/(2a)$ (o observant que el vèrtex sempre està al punt mig de les arrels). Per a la justificació: qualsevol paràbola amb arrels a 0 i a 5 s'ha de poder escriure com $f(x) = ax(x - 5)$, per a algun valor de $a$. Imposa que aquesta expressió, avaluada al punt mig de les arrels, doni exactament $6{,}25$: obtindràs una equació que et permetrà trobar el valor de $a$ de manera única.}

\end{document}
